\section*{第 \ref{ch:g6:shapes} 章——三角形与四边形}

\begin{solution}{exo:g6:shapes:1}
按作图方法：作\omterm{def:g6:lines:objects}{线段} $[AB]$，长 $6$ cm；画\omterm{def:g3:shapes:shapes}{半径}为 $5$ cm、
圆心为 $A$ 的弧，再画\omterm{def:g3:shapes:shapes}{半径}为 $4$ cm、圆心为 $B$ 的弧；
两弧的交点即为 $C$。
\end{solution}

\begin{solution}{exo:g6:shapes:2}
底边 $[EF]$ 长 $3$ cm，画\omterm{def:g3:shapes:shapes}{半径}为 $5$ cm、
圆心分别为 $E$ 和 $F$ 的两段弧，交点为 $D$。两底角 $\widehat{DEF}$
与 $\widehat{DFE}$ 相等（各约 $72.5^\circ$）——\omterm{def:g6:shapes:triangles}{等腰三角形}的对称性
（见\cref{ex:g6:symmetry:proves}）。
\end{solution}

\begin{solution}{exo:g6:shapes:3}
作长 $4$ cm 的\omterm{def:g6:lines:objects}{线段}，再画两段\omterm{def:g3:shapes:shapes}{半径}为 $4$ cm 的弧，圆心分别在
\omterm{def:g6:lines:objects}{线段}的两个端点。每个角都等于 $60^\circ$。
\end{solution}

\begin{solution}{exo:g6:shapes:4}
(a) \omterm{def:g6:shapes:triangles}{等边三角形}（因此也是\omterm{def:g6:shapes:triangles}{等腰三角形}）。\quad
(b) \omterm{def:g6:shapes:triangles}{等腰三角形}（$5 = 5$）。\quad
(c) \omterm{def:g3:shapes:rightangle}{直角}\omterm{def:g6:shapes:triangles}{等腰三角形}（两个 $45^\circ$ 角
所对的两边相等）。\quad
(d) \omterm{def:g6:shapes:triangles}{等边三角形}（三个角相等）。
\end{solution}

\begin{solution}{exo:g6:shapes:5}
斜边是与\omterm{def:g3:shapes:rightangle}{直角}相对的边，总是最长的
边。在 $MNP$ 中，\omterm{def:g3:shapes:rightangle}{直角}位于 $N$，斜边是 $[MP]$：
它不经过\omterm{def:g3:shapes:rightangle}{直角}的\omterm{def:g5:solids:def}{顶点} $N$。
\end{solution}

\begin{solution}{exo:g6:shapes:6}
\emph{1. 正确}：正方形有四个\omterm{def:g3:shapes:rightangle}{直角}，这正是
长方形的定义。

\emph{2. 错误}：一个 $5 \times 3$ 的长方形两边不相等。

\emph{3. 正确}：正方形四条边相等，这正是
\omterm{def:g6:shapes:quadrilaterals}{菱形}的定义。

\emph{4. 正确}：见\cref{ex:g6:shapes:recognize}。
\end{solution}

\begin{solution}{exo:g6:shapes:7}
用垂线作图，方法同\cref{ex:g6:shapes:program}。
两条对角线长约 $5.8$ cm：长方形的对角线
长度相等（见\cref{prop:g6:shapes:diagonals}）。
\end{solution}

\begin{solution}{exo:g6:shapes:8}
画长 $6$ cm 的\omterm{def:g6:lines:objects}{线段}，取中点 $O$；过 $O$ 作垂线；
在垂线上各取距离 $2$ cm 的一点，分居 $O$ 两侧。依次连接
四个端点即得\omterm{def:g6:shapes:quadrilaterals}{菱形}（四个\omterm{def:g5:solids:def}{顶点}的距离分别为
$3$ cm 和 $2$ cm，均沿垂线以 $O$ 为出发点，
故四边因对称而相等）。
\end{solution}

\begin{solution}{exo:g6:shapes:9}
一种正确的作图程序：
1. 作\omterm{def:g6:lines:objects}{线段} $[BC]$，使 $BC = 4$ cm；
2. 作圆心为 $B$、\omterm{def:g3:shapes:shapes}{半径}为 $6$ cm 的弧，位于 $[BC]$ 上方；
3. 作圆心为 $C$、\omterm{def:g3:shapes:shapes}{半径}为 $6$ cm 的弧，位于 $[BC]$ 上方；
4. 记两弧的交点为 $A$；
5. 连接 $[AB]$ 和 $[AC]$。
\end{solution}

\begin{solution}{exo:g6:shapes:10}
无论 $M$ 在圆上的位置如何，角 $\widehat{AMB}$
都为 $90^\circ$：三角形 $AMB$ 的\omterm{def:g3:shapes:rightangle}{直角}位于 $M$
只要 $[AB]$ 为\omterm{def:g4:geometry:circle}{直径}。证明留待
\cref{ch:g8:cosine}。
\end{solution}

\begin{solution}{exo:g6:shapes:11}
对角线互相平分且长度相等：
这正是\emph{长方形}的性质
（见\cref{prop:g6:shapes:diagonals}）。并不要求对角线
互相\omterm{def:g6:lines:perp}{垂直}，因此不必是\omterm{def:g6:shapes:quadrilaterals}{菱形}或正方形——长方形
才是正确答案（正方形也满足条件，但并非
必然如此）。
\end{solution}

\begin{solution}{exo:g6:shapes:12}
$4 \times 4$ 的点阵中共有 $3 \times 3$ 个方格。
沿格线的正方形：边长取 $1$、$2$ 或 $3$，共有 $9 + 4 + 1 = 14$
个正方形。另有倾斜的正方形：较小的一类，其边为
$1 \times 1$ 方格的对角线（如 $(1,0)$、$(2,1)$、$(1,2)$、
$(0,1)$）——共四个；另有较大的一类，其边为
$2 \times 1$ 长方形的对角线（如 $(1,0)$、$(3,1)$、$(2,3)$、
$(0,2)$）——共两个。此类正方形的四条边都相等，
因为它们都是全等长方形的对角线。合计：
$14 + 4 + 2 = 20$ 个正方形。（只要倾斜例子正确，
沿格线的 $14$ 个也算对，即为本阶段的成功。）
\end{solution}

\begin{solution}{pb:g6:shapes:1}
\textbf{1.} 四条边都等长（约
$4.2$ cm），四个角都是\omterm{def:g3:shapes:rightangle}{直角}：图形是一个
\emph{正方形}。（对角线相等得\omterm{def:g3:shapes:rightangle}{直角}，与长方形中
相同；互相\omterm{def:g6:lines:perp}{垂直}的对角线得四边相等，与\omterm{def:g6:shapes:quadrilaterals}{菱形}中
相同；两者兼有即得正方形，
由\cref{prop:g6:shapes:diagonals}逆向得出。）

\textbf{2.} 图形是\emph{\omterm{def:g6:shapes:quadrilaterals}{菱形}}：四条边相等（约 $4.7$ cm），
但没有\omterm{def:g3:shapes:rightangle}{直角}——正是
\cref{exo:g6:shapes:8}的作图。

\textbf{3.} 四个角都是\omterm{def:g3:shapes:rightangle}{直角}，但四条边分成
两种不同的长度：这是一个\emph{长方形}，
即\cref{exo:g6:shapes:11}所示的情形。

\textbf{4.} 测量之后，没有任何等长的结论成立，也没有
\omterm{def:g3:shapes:rightangle}{直角}出现；但对边两两相等，
看上去也互相\omterm{def:g6:lines:perp}{平行}。这是一般的“倾斜的长方形”，它的
对角线只是互相平分——即
\emph{\omterm{def:g6:lines:perp}{平行}\omterm{def:g4:geometry:polygon}{四边形}}，在\cref{ch:g7:central}中命名并研究。

\textbf{5.} 对角线相交于它们的公共中点：
\begin{center}
\renewcommand{\arraystretch}{1.2}
\begin{tabular}{l|l}
对角线\dots & 形状 \\
\hline
相等且\omterm{def:g6:lines:perp}{垂直} & 正方形 \\
只\omterm{def:g6:lines:perp}{垂直} & \omterm{def:g6:shapes:quadrilaterals}{菱形} \\
只相等 & 长方形 \\
两者都不满足 & 未命名的一类（见\cref{ch:g7:central}） \\
\end{tabular}
\end{center}

\textbf{6.} 每个角上的三角形沿每条对角线各有一条直角边：
即 $8$ cm 与 $5$ cm 的\omterm{def:g3:division:half}{一半}，即 $4$ cm
和 $2.5$ cm，两者之间夹\omterm{def:g3:shapes:rightangle}{直角}——四个三角形的数据都相同。
由相同的两条直角边及其间的\omterm{def:g3:shapes:rightangle}{直角}作出的三角形，
彼此全等。
（作图时圆规的开度完全相同）。它们的四条斜边——即图形的边——
因此都相等：
图形是\omterm{def:g6:shapes:quadrilaterals}{菱形}。

\textbf{7.} 把圆规张开到一条边的长度，然后
沿图形绕行一周：若针尖与笔尖恰好落在
每对相邻\omterm{def:g5:solids:def}{顶点}上，则四条边都与这同一开度相等。
无需直尺。

\textbf{8.} 作图步骤：
\begin{enumerate}
  \item 画\omterm{def:g6:lines:objects}{线段} $[AC]$，使 $AC = 6$ cm，并
        取中点 $O$（$3$ cm 为其到 $A$ 的距离）。
  \item 作 $(AC)$ 的垂线，使其过 $O$。
  \item 在垂线上取点 $B$、$D$，使它们到中点的距离均为 $3$ cm，
        且分居 $O$ 的两侧。
  \item 依次连接 $A$、$B$、$C$、$D$。
\end{enumerate}
对角线相等（$6$ cm）、互相\omterm{def:g6:lines:perp}{垂直}、互相平分：
对照第 5 题的表，$ABCD$ 是正方形。

\textbf{9.} 测量约得 $4.2$ cm，\emph{比}
$3$ cm 更长。常见的猜测“半条对角线”是错误的。
可证明的部分是：从\omterm{def:g5:solids:def}{顶点} $A$ 走到下一\omterm{def:g5:solids:def}{顶点} $B$，
可以沿边 $[AB]$ 直接走，也可以绕道经过中心
$A \to O \to B$，共 $3 + 3 = 6$ cm；
沿直线的那条边更短，故边长小于 $6$ cm。
它也\emph{大于} $3$ cm，这一点由直尺可以看出；
在\cref{ch:g8:pythagoras}中将证明这一点，并给出精确值
（$3$ cm 乘以 $2$ 的平方根，约为 $4.24$ cm）。

\textbf{10.} 在 $H$ 处的四个角三角形，现在分成
\emph{两对}相配的三角形：其中两个直角边为 $3$ cm 和 $2.5$ cm
（斜边为 $[AB]$ 与 $[AD]$），另两个直角边为 $5$ cm 和
$2.5$ cm（斜边为 $[CB]$ 与 $[CD]$）。于是 $AB = AD$ 且
$CB = CD$：\omterm{pb:g6:shapes:1}{筝形}有两对\emph{相邻}的边分别相等
（\omterm{def:g6:shapes:quadrilaterals}{菱形}则是四边全都相等的特殊\omterm{pb:g6:shapes:1}{筝形}）。

\textbf{11.} \omterm{def:g4:geometry:polygon}{四边形}：$2$ 条对角线。\omterm{def:g4:geometry:polygon}{五边形}：$5$ 条。
\omterm{def:g4:geometry:polygon}{六边形}：$9$ 条。

\textbf{12.} 从\omterm{def:g4:geometry:polygon}{六边形}的一个\omterm{def:g5:solids:def}{顶点}出发，对角线连至
除自身及两个相邻\omterm{def:g5:solids:def}{顶点}外的每个\omterm{def:g5:solids:def}{顶点}：共 $6 - 3 = 3$ 条。
六个\omterm{def:g5:solids:def}{顶点}共引出 $6 \times 3 = 18$ 个对角线端点，
——但每条对角线在\emph{两端}各被计数一次，
即从两个端点各计一次，恰如握手问题中被两位握手者各计一次
（见\cref{pb:g6:wholes:1}）。真正的条数为：
$18 \div 2 = 9$，与图形相符。

\textbf{13.} 十个\omterm{def:g5:solids:def}{顶点}：每个\omterm{def:g5:solids:def}{顶点}引出 $10 - 3 = 7$ 条，
于是共 $10 \times 7 \div 2 = 35$ 条对角线。二十个\omterm{def:g5:solids:def}{顶点}：
共 $20 \times 17 \div 2 = 170$ 条对角线。

\textbf{14.} 试算候选值：$10$ 个\omterm{def:g5:solids:def}{顶点}得 $35$ 条（第 13 题），
$11$ 个\omterm{def:g5:solids:def}{顶点}得 $11 \times 8 \div 2 = 44$ 条。找到：
该\omterm{def:g4:geometry:polygon}{多边形}有 $11$ 个\omterm{def:g5:solids:def}{顶点}。

\textbf{15.} 三角形根本没有对角线：每个\omterm{def:g5:solids:def}{顶点}的
“非相邻\omterm{def:g5:solids:def}{顶点}”都是\,\dots\ 没有，因为另外两个\omterm{def:g5:solids:def}{顶点}
都是它的相邻\omterm{def:g5:solids:def}{顶点}。方法的结果一致：
每个\omterm{def:g5:solids:def}{顶点}各得 $3 - 3 = 0$ 条，共 $0$ 条。握手问题与对角线
其实是同一种计数，只是换了装束：从一组中选出的每一对，
只计一次——数学总是反复利用
它的好想法。
\end{solution}
